\section{Problema de la mochila (Knapsack)}

\subsection{Introducción}

Dentro del área de problemas de Cargas y Empaquetamientos, los más sencillos son los problemas de una carga, conocidos
como problemas de la mochila (knapsack) según la terminología de Dantrzig. En estos problemas se trata de minimizar un coste
o maximizar un beneficio al seleccionar una serie de elementos para ``introducir en la mochila'', cumpliendo restricciones
sobre las características de estos mismos (peso por ejemplo).

\subsection{Problema de la Mochila 0-1}

De forma general consideraremos

\begin{center}
    \begin{tabular}{rcl}
        $n$   & = & número elementos para elegir \\
        $p_j$ & = & el valor del objeto $j$ \\
        $w_j$ & = & el peso del objeto $j$ \\
        $c$   & = & peso total asumible
    \end{tabular}
\end{center}

\noindent Por tanto el problema de la mochila binaria o 0-1 se puede formular de la siguiente manera

\begin{alignat}{3}
    \text{maximizar:}   &   \sum_{j=1}^n p_jx_j       \\
    \text{sujeto a:}    &   \sum_{j=1}^n w_jx_j\leq c \\
                        &   x_j \in {0, 1} \quad j=1,\dots,n
\end{alignat}

\noindent donde $x_j$ es la variable binaria que indica si el elemento $j$-ésimo está presente o no en la solución

Existe una versión entera de este problema conocida como \textit{Problema de la Mochila Acotado}, donde $x_j$ indica la
cantidad de dicho elemento que se va a incluir en vez de la presencia en si misma, por lo que $n$ ahora no es el número
total de elementos sino el número de elementos diferentes. Se definen también cotas para la cantidad de cada elemento $u_j$
tal que $0\leq x_j\leq u_j$.

En ocasiones (y como ya se ha mencionado) se puede considerar el problema inverso, el de minimización, donde tratamos de 
minimizar la función objetivo. Formular este problema es tan sencillo como invertir la restricción tal que

\begin{alignat}{3}
    \text{minimizar:}   &   \sum_{j=1}^n p_jx_j       \\
    \text{sujeto a:}    &   \sum_{j=1}^n a_jx_j\geq b \\
                        &   x_j \in {0, 1} \quad j=1,\dots,n
\end{alignat}

\noindent donde en esta ocasión $a_j$ es una caracteristica medida de los elementos y $b$ la cota mínima total.

\newpage

\subsubsection{Algoritmos Greedy}

Diseñar un algoritmo voráz para un problema mochila 0-1 es muy sencillo, basta con elegir una función o criterio ``greedy''
$g_j$ e ir introduciendo en la mochila los elementos que maximicen dicho criterio de forma sucesiva hasta que no entren
más (es decir mientras se satisfaga la condición). 

\vspace*{2mm}

La primera opción es considerar el valor de los elementos como criterio del algoritmo, el problema que tiene esto es que no 
estamos considerando el peso de los objetos. Del mismo modo, considerar el peso de los objetos como el criterio tampoco es lo 
suficientemente completo, pues no sabemos nada respecto del valor de estos por lo que no aporta información sobre nada 
relacionado con la función objetivo. 

Una posible alternativa es elegir el criterio como el cociente del valor por el peso, de forma que $g_j=\frac{p_j}{w_j}$,
lo cual parece lógico pues establecemos que lo valioso que es un objeto para nosotros depende de la relación entre el valor
y el peso. Este último criterio es el más utilizado, pero no significa que se encuentre la solución óptima a partir de él.

\vspace*{2mm}

\begin{ejemplo}
    En este ejemplo se muestra en una tabla una serie de elementos con sus correspondientes valores, pesos y cocientes
    
    \begin{center}
        \begin{tabular}{c|cccccccc}
            \hline
            elemento & 1 & 2 & 3 & 4 & 5 & 6 & 7 & 8 \\ \hline
            valor & 15 & 100 & 90 & 60 & 40 & 15 & 10 & 1 \\
            peso & 2 & 20 & 20 & 30 & 40 & 30 & 60 & 10 \\ \hline
            $\frac{\text{valor}}{\text{peso}}$ & 7.5 & 5 & 4.5 & 2 & 1 & 0.5 & 0.16 & 0.1 \\
            \hline
        \end{tabular}
    \end{center}

    Sabiendo que el peso máximo es 102 $c=102$, se obtienen los siguientes resultados (representados en un array de
    índices, no con las variables binarias, ordenados por elección)

    \begin{itemize}
        \item Tomando como criterio el valor, el array resultante es $[2, 3, 4, 1, 6]$ con un valor total de 280.
        \item Tomando como criterio el peso, el array resultante es $[7, 5, 1]$ con un valor total de 65.
        \item Por último, tomando como criterio el cociente, el array resultante es $[1, 2, 3, 4, 6]$ con un valor de 280.
    \end{itemize}
\end{ejemplo}

\subsubsection{Relajación Lineal}

La relajación lineal de un Problema de la Mochila 0-1 consiste en eliminar la restricción de variables binarias y sustituirla
por una cota inferior o superior para $x_j$ tal que 

\begin{alignat}{3}
    \text{maximizar:}   &   \sum_{j=1}^n p_jx_j       \\
    \text{sujeto a:}    &   \sum_{j=1}^n w_jx_j\leq c \\
                        &   0\leq x_j\leq 1 \quad j=1,\dots,n
\end{alignat}

\noindent es decir, que permitimos a un elemento ser elegido de forma parcial.

\newpage

A la solución factible obtenida mediante este procedimiento la denotaremos $z_{LKP}$. Si redondeamos $z_{LKP}$ (truncar)
obtendremos otra solución factible a la que llamaremos $z'$. Se verifica que
\[
    z' \leq z \leq z_{LKP}
\]
donde $z$ es la solución óptima del problema de la mochila. Por lo que estas dos soluciones factibles son una cota inferior
y superior de la solución óptima.

\subsubsection{Mejoras de las soluciones greedy}

Para mejorar la solución que se obtiene mediante el algoritmo greedy puede usarse cualquier método que combine la 
construcción greedy con la aleatorización en la forma de un método multistart, como los métodos random greedy y greedy 
aleatorizado.

Otra opción es optimizar un resultado mediante busqueda local. El procedimiento consiste en, dada una solución $x$ a un
problema de la mochila, se selecciona un entorno de soluciones cercanas y se compara la solución obtenida contra la que
máxima dentro del entorno. Si la solución del entorno mejora la actual se selecciona y se vuelve a aplicar el procedimiento,
extrayendo un nuevo entorno y seleccionando la mejor hasta que ninguna supere la previamente seleccionada. La idea detras
de este procedimiento es encontrar un máximo local de la función solución. Al tratarse de un espacio discreto debemos definir
lo que es un entorno de la solución. Este se obtiene mediante intercambios de variables en el array solución, se eligen 
$p=1,2,\dots$ elementos en el array solución que se deben intercambiar por otros $q=1,2,\dots$ que no esten presentes y que
verifiquen que formen una solución factible. Cada una de estas permutaciones se trata de una solución del entorno de $x$.

\subsection{Problema de la Mochila Múltiple}

En el Problema de la Mochila Múltiple se tienen $m$ mochilas o contenedores. Para cada $i=1,\dots,m$ $c_i$ es la capacidad
de la mochila $i$-ésima. El modelo completo es el siguiente

\begin{alignat}{3}
    \text{maximizar:}   &   \sum_{j=1}^n p_jx_j                             \\
    \text{sujeto a:}    &   \sum_{j=1}^n w_jx_j\leq c_i \quad i=1,\dots,m   \\
                        &   \sum_{i=1}^m x_{ij}\leq 1 \label{eq:62}           \\
                        &   x_j \in {0, 1} \quad j=1,\dots,n
\end{alignat}

\noindent donde aparecen $mn$ variables binarias nuevas, $x_{ij}$, que representan la presencia o no del elemento $j$ en la 
mochila $i$. La restricción \ref{eq:62} hace referencia al hecho de que un elemento solo puede estar presente en una mochila
y no en varias a la vez.

\subsubsection{Algoritmos greedy}

Una posible implementación sencilla de un algoritmo greedy para el Problema de la Mochila Múltiple es tomar la mochila de
mayor tamaño y ``llenarla'' usando el criterio del cociente antes descrito. Una vez ``llena'' se pasa a la siguiente mochila
más grande y se repite el procedimiento hasta no poder incluir ningun elemento más en la última mochila. Aplicar busqueda 
local basada en intercambios sobre la solución greedy suele proporcionar buenos resultados.

\subsection{Problema de la Mochila Multidimensional}

En esta ocasión se tienen $m$ caracteristicas sobre los elementos y $m$ restricciones sobre dichas características. Llamaremos
$a_{ij}$ al valor de la caracteristica $i=1,\dots,m$ del objeto $j$ y $b_i$ al límite sobre dicha característica en el
problema. La formulación del problema es la siguiente

\begin{alignat}{3}
    \text{maximizar:}   &   \sum_{j=1}^n p_jx_j                                 \\
    \text{sujeto a:}    &   \sum_{j=1}^n a_{ij}x_j\leq b_i \quad i=1,\dots,m    \\
                        &   x_j \in {0, 1} \quad j=1,\dots,n
\end{alignat}

Este modelo también se aplica sobre problemas de inversiones multiperíodo o a la selección de proyectos. Aquí $n$ es el 
número de posibles inversiones o proyectos, $m$ el número de períodos o recursos, $b_i$ es el capital (o recurso) disponible
en el periodo $i$, $a_{ij}$ la inversión (o cantidad de recurso) requerida por el proyecto $j$ en el periodo $i$ y $p_j$ la
rentabilidad o tasa de ganancia esperada para el proyecto o inversión.

\subsubsection{Algoritmos greedy}

Igual que en el caso de la mochila simple, para construir una solución iterativamente mediante un algoritmo greedy hay que
asignar un coeficiente o criterio greedy a cada elemento/proyecto. Para el caso multidimensional se han propuesto diferentes
criterios para definir los coeficientes. Se destacan versiones de $g_j=\frac{p_j}{\sum_{i=1}^m \mu_ia_{ij}}$ donde $\mu_i$ 
indica un valor dado al recurso $i$ en función de su escasez.

\begin{itemize}
    \item Mediante el cociente entre el beneficio y la suma agregada de recursos que consume
    \[
        g_j=\frac{p_j}{\sum_{i=1}^m a_{ij}}
    \]
    Equivale a tomar $\mu_i=1$, es decir, tratar todos los recursos por igual
    \item Teniendo en cuenta la escasez de los recursos
    \[
        g_j=\frac{p_j}{\sum_{i=1}^m \frac{a_{ij}}{b_i}}
    \]
    En el caso de subastas, a este coeficiente se le conoce como el precio escalado normalizado de la puja u oferta $j$.
    \item Seleccionando $\mu_i$ como el precio sombra de la restricciín $i$-ésima en la solución óptima de relajación lineal.
\end{itemize}

\noindent También se tienen en cuenta coeficientes basados en la solución del problema lineal. A cada elemento se le asigna
como coeficiente greedy la solución de la relajación lineal, es decir, $g_j=x_{j}.sol$.

\newpage

\noindent En algunas versiones del algoritmo greedy, como el algoritmo de London u Michaelis, el criterio greedy se define y 
calcula en cada iteración en función de los recursos usados hasta ese momento.

\vspace*{5mm}

Al igual que para el resto de Problemas de la Mochila, se pueden aplicar mecanismos de busqueda local mediante intercambios
para encontrar máximos locales que mejoren la solución greedy y mecanismos de aleatorización en los elementos y en los 
intercambios.